\documentclass[UTF8,a4paper,10pt,fontset=fandol]{ctexart}
\usepackage[margin=20mm]{geometry}
\usepackage{amsmath,amssymb,amsthm,booktabs,array,tabularx,graphicx,pgf}
\usepackage{enumitem,microtype,caption,hyperref}
\setsansfont{TeX Gyre Heros}
\hypersetup{colorlinks=true,linkcolor=blue!45!black,citecolor=blue!45!black,urlcolor=blue!45!black}
\captionsetup{font=small,labelfont=bf}
\setlist{nosep,leftmargin=1.6em}
\setlength{\emergencystretch}{2em}
\newtheorem{proposition}{命题}
\newsavebox{\pcnfigurebox}
\newcommand{\pcn}{\mathcal D_q}
\newcommand{\todo}[1]{\textbf{TODO：#1}}
\title{路径一致负例：单路径结构化检索中的监督冲突与修正}
\author{研究稿；作者信息待填写}
\date{2026年10月4日}
\begin{document}
\maketitle

\begin{abstract}
路径弱监督可以减少结构化检索对逐转移标注的依赖，但选择一条答案路径并不能使其余实际采样转移成为可靠负例。我们将实际负采样集合中同时属于另一条完整最短答案路径的转移定义为路径一致负例（Path-Consistent Negative, PCN），以完整 prefix--transition--suffix 条件识别冲突，再保持标签、撤销相应负损失。该方法不增加参数，也不改变候选生成和推理。在 WikiTopics、MetaQA-3hop、PathQuestion-2H3H 及 KQAPro 实体单主题子集上，PCN 占实际负例的比例为 4.59\%、11.81\%、8.80\% 和 1.45\%。共享五个种子的实验中，屏蔽相对原始监督的 APC-MRR 分别提高 0.257、1.323、2.179 和 1.132 个百分点，且均优于逐题等量随机屏蔽。正例重标记未呈现统一劣势；另行重建全最短路正监督在两个新基准上均高于屏蔽。结果将最小修正的价值限定为保留既有监督集合时的冲突处理，而非最佳多路径训练方案。完整路径身份可揭示结构监督的不一致，但不能独自提供语义正例标注。
\end{abstract}
\noindent\textbf{关键词：}结构化检索；路径弱监督；负采样；监督一致性；知识图谱

\section{引言}
结构化知识检索需要从图中的大量事实中选出能够共同支撑答案的关系链。与逐条标注相关性相比，利用主题实体和答案之间的路径产生监督更容易扩展，知识图谱问答与超图检索均采用了路径或关系链作为学习对象\cite{metaqa,pq,hyperrag}。这种监督把答案层面的标注转化为局部转移标签，其可靠性取决于路径选择与负采样是否使用一致的依据。

单路径构造通常对每个答案选择一条最短路径，将所选转移记为正例，再从附近候选中抽取负例。这里存在一个隐含推断：候选未出现在所选正路径中，便可以获得确定负标签。当答案路径唯一时，这种选择不会遗漏其他等长答案链；存在多条完整最短路径时，未被选择的路径却可能提供同样的结构依据。单条路径的选择结果因此不能充分支持其他候选的负标签。

图~\ref{fig:overview} 给出这一结构监督冲突（structural supervision conflict）。主题 $s$ 经 $x$ 或 $y$ 都能以两跳到达答案 $a$，监督只选择 $s\to x\to a$。若 $y\to a$ 被实际负采样器抽中，它在损失中受到负监督，同时又具有完整最短答案路径成员资格。仅有替代路径并不构成本文研究的 PCN；该转移必须真实出现在本次训练的负例集合中。另一方面，绕远前缀上的 $z\to a$ 虽然局部朝答案前进，却不属于最短完整链。

这种冲突在多个数据来源中实际出现。五种子训练样本中，WikiTopics 的 9,885,005 个负例包含 453,481 个 PCN，MetaQA 的 14,162,090 个负例包含 1,672,257 个 PCN。新增的 PathQuestion 和 KQAPro 实体子集也分别包含 5,048 和 546 个 PCN。这些计数来自采样后张量，而不是把图上全部替代边视为训练噪声。跨图结构的现象使得监督构造本身值得检验，而不只是讨论某个领域的检索得分。

我们的最小一致性修正（minimal consistency correction）只撤销 PCN 的负损失：标签仍为零，权重变为零。这一操作称为路径一致负例屏蔽（Path-Consistent Negative Masking, PCN Mask）。选择来自证据的不对称性：完整最短答案路径足以反驳“因未被选择而确定为负”的推导，但问题语义可能排除某些结构捷径，路径成员资格不能自动支持确定正标签。因此，我们同时比较原始监督、等量随机屏蔽、PCN 屏蔽和 PCN 正例重标记，并另行考察改变正例构造的全最短路监督。

四组数据的宏观结果支持 PCN 屏蔽改善 APC-MRR，且其差值高于逐题等量随机屏蔽。证据同时限定了这一结论：屏蔽与重标记没有统一优劣，全最短路构造在两个新基准上更高，部分领域及 Reach 指标也没有改善。本文贡献依次为：识别实际采样中的结构监督冲突；给出完整最短路径条件及跨图实现；提供保留评分器和推理的最小监督撤销；以四组数据、配对控制和负结果说明何种结论成立、何种结论尚不成立。

\begin{figure}[tbp]
\centering
\begin{lrbox}{\pcnfigurebox}
@@overview_pgf@@
\end{lrbox}\resizebox{\linewidth}{!}{\usebox{\pcnfigurebox}}
\caption{完整路径身份与监督选择之间的冲突。左：两条完整两跳答案路径、所选路径和实际采样的 PCN；灰色三跳链展示局部靠近答案但 prefix 绕远的反例。右：四策略共享候选与评分模型，只改变训练时的标签或权重。图为解释定义的构造示例，不是数据集中的计数案例。推理只接收问题、topic 和候选特征。}\label{fig:overview}
\end{figure}

\section{单路径弱监督中的结构冲突}
\subsection{图、转移与监督集合}
设有问题 $q$、主题实体 $s_q$、答案集合 $A_q$ 和固定有向图 $G$。普通知识图谱的转移为 $t=(u,r,v)$，其图距离代价 $c(t)=1$。对于 WikiTopics，检索图采用实体与超边的二部表示，转移 $t=(u,e,v)$ 表示两条二部边，距离代价为 $c(t)=2$。逆向转移在固定图中显式构造，不能只因尾实体相同就忽略方向或关系身份。

对每个训练答案 $a$，记所有允许距离范围内的最短 $s_q$--$a$ 路径集合为 $\mathcal P^*(s_q,a)$。构造器从中稳定选择一条 $\widehat P_{q,a}$，所选正例集合及完整最短路径并集分别为
\begin{align}
S_q &= \bigcup_{a\in A_q:\,\mathcal P^*(s_q,a)\ne\varnothing} T(\widehat P_{q,a}),\\
U_q &= \bigcup_{a\in A_q} \bigcup_{P\in\mathcal P^*(s_q,a)}T(P).
\end{align}
“单路径”指每个答案选一条路径；多答案问题的正例可能是多条所选路径的并集。$N_q$ 是该问题、该种子下\emph{实际}被负采样器选出的转移，采样排除 $S_q$。不同图适配保留既有的路径引导候选构造及逆向正路径排除规则。本轮普通 KG 实际负例数为正例数与可用负例池大小的较小者；WikiTopics 保留有限次采样尝试，允许不足额，不把预期数量充当观测数量。

\subsection{冲突定义与解释范围}
路径一致负例定义为
\begin{equation}
\pcn=N_q\cap U_q.
\label{eq:pcn}
\end{equation}
因为 $N_q$ 排除了所选正例，PCN 在当前构造中具有 $y_t=0$，却属于完整最短答案路径。它不是图中的任意“可能有用边”，也不包括从未被采样的替代路径转移。

这一身份指出的是结构监督规则的内部冲突：最短答案路径资格既然被用来支持所选转移的正监督，仅凭选择器遗漏就不能支持相同资格转移的确定负监督。它并不是关于自然语言语义的矛盾定理。无标注的结构路径可以经过错误关系、语义捷径或与问题约束无关的事实；本文不将所有 PCN 都称为语义正确的证据。

\section{通过完整路径识别 PCN}
\subsection{完整距离条件}
令 $d(s,u)$ 为固定图中的有向最短距离。对已知属于图的有效转移 $t=(u,\cdot,v)$，其完整最短答案路径成员条件为
\begin{equation}
\exists a\in A_q:\quad d(s_q,u)+c(t)+d(v,a)=d(s_q,a)<\infty.
\label{eq:distance}
\end{equation}
距离按普通 KG 或二部超图使用一致单位，答案限定于训练构造器支持的三跳范围。只有满足式~\eqref{eq:distance} 且属于 $N_q$ 的转移才进入 $\pcn$。条件检查从 topic 到转移的完整 prefix、转移本身和到答案的 suffix；单独检查 $d(v,a)<d(u,a)$ 会误纳已经绕远的前缀。

\begin{proposition}
对于固定有向图中有效、正代价的转移，式~\eqref{eq:distance} 当且仅当该转移属于至少一条完整最短 $s_q$--$a$ 路径。
\end{proposition}
\begin{proof}
若转移位于一条最短路径上，其 prefix 与 suffix 必须分别最短，否则替换后会得到更短路径，故距离相加等于 $d(s_q,a)$。反之，连接一条最短 prefix、该有效转移及一条最短 suffix，长度恰为 $d(s_q,a)$，构成最短答案链。正代价排除相交拼接产生可删循环而仍达到最短长度的情况。
\end{proof}

命题证明完整\emph{图路径}身份，不证明路径满足问题的全部语义约束。它也没有宣称最短链是唯一有效证据；更长但语义正确的路径不在当前定义内。

\subsection{答案最短路径 DAG 的实现}
直接为每个候选和答案重复计算 suffix 距离并无必要。先对 topic 做有界 BFS 得到 $d_s$ 和所有最短前驱，再从范围内可达答案沿前驱回溯，得到完整答案最短路径节点并集 $V_q^*$。在有效转移上，式~\eqref{eq:distance} 等价于
\begin{equation}
d_s(v)=d_s(u)+c(t),\qquad v\in V_q^*.
\label{eq:dag}
\end{equation}
第二项保证存在通往某个答案的最短 suffix，第一项保证 prefix 没有绕远。严格增加的距离使该结构为 DAG。普通 KG 用代价 1，超图二部转移用代价 2，调用同一个判定函数；关系/超边的有效性由图适配器确认。

实现缓存同一 topic 的 BFS，按问题答案做回溯，不枚举每一条完整路径。一次 topic 遍历成本为 $O(|V_h|+|E_h|)$，其中 $h$ 是三跳范围；答案回溯与实际负例检查分别与答案 DAG 和 $|N_q|$ 成正比。路径多重性用 DAG 上的动态计数获得，路径数较大时仍不展开路径列表。该分析没有额外学习器，也不需要模型置信度阈值。

\section{最小一致性修正与对照}
\subsection{共享评分器与损失}
评分器保持 HyperRetriever 的发布接口\cite{hyperrag}。冻结 GTE-large-en-v1.5 提供问题、head、关系或超边文本、tail 的四个 1,024 维向量，加上只由 topic 和候选结构计算的 30 维 DDE。4,126 维特征进入 256 维 ReLU 隐层与单输出层。监督策略不修改这个函数 $f_\theta(q,t)$。

对一个 batch $B$，加权二元交叉熵为
\begin{equation}
\mathcal L_B=\frac{\sum_{i\in B}w_i\,\ell_{\mathrm{BCE}}(f_\theta(q_i,t_i),y_i)}{\sum_{i\in B}w_i}.
\label{eq:loss}
\end{equation}
PCN Mask 保持 $y_i$，令 $t_i\in\pcn$ 时 $w_i=0$，其余 $w_i=1$。对应负项的梯度贡献为零，其余样本保留。归一化分母随有效权重改变，因此不声称其他样本的归一化梯度完全不变。实现跳过有效权重和为零的 batch；没有为 PCN 生成额外正例，也不增加参数。

\subsection{四种监督策略}
表~\ref{tab:strategies} 将策略差别限定在共享训练集合上。Random Mask 对每个问题，从 $N_q$ 无放回随机屏蔽恰好 $|\pcn|$ 个候选，可随机选到 PCN；不根据路径身份排除它们。它控制减少负监督的数量，并保留随机选择的不确定性。Positive Relabel 只对真实 PCN 改标，不扩大采样集合。

\begin{table}[tbp]\centering
\caption{四策略在同一实际负例集合上的操作。所选正例及其权重保持不变。}\label{tab:strategies}
\begin{tabular}{lccc}\toprule
策略 & PCN 标签 & 屏蔽位置 & 权重 \\\midrule
Original & 0 & 无 & 全 1 \\
Random Mask & 0 & $R_q\subseteq N_q,\ |R_q|=|\pcn|$ & $R_q$ 为 0 \\
PCN Mask & 0 & $\pcn$ & $\pcn$ 为 0 \\
Positive Relabel & 1 & 无 & 全 1 \\\bottomrule
\end{tabular}
\end{table}

全最短路径监督是另一个问题：先把 $S_q$ 换为 $U_q$，再构造候选负例池和采样。它改变训练候选及正负数量，不能与只改 $\pcn$ 标签的 Positive Relabel 混同。我们把它作为普通 KG 上的额外构造对照，单独报告，四策略主比较始终共享原始训练候选。

\subsection{训练与推理解耦}
PCN 识别只使用训练答案。evaluation 候选从 topic 出发，由固定语义相似度和图规则产生，宽度 32、至多三跳。普通 KG 适配对当前 frontier 的转移按问题与 head/关系/tail 的平均余弦相似度选择；WikiTopics 使用既有超图候选流程。排序时仅输入特征和训练后评分器，不查询 PCN 集合或答案。测试答案只在候选准备之后用于完整路径指标、Oracle 和说明性结构分析。

\section{跨数据集实验设计}
\subsection{研究问题与数据范围}
实验依次回答：RQ1，实际 PCN 是否跨数据来源存在；RQ2，定向撤销是否改善排序；RQ3，改善能否由等量减少负监督解释；RQ4，屏蔽与正例重标记如何比较；RQ5，图结构和问题分布如何对应收益。额外 RQ6 比较监督构造阶段的全最短路并集。

WikiTopics 包含十一领域，采用 QID 保持实体身份、冻结文本标签和 entity--hyperedge 二部结构\cite{hyperrag}。复用的固定图包含发布的训练图与推理图，是已知语料的传导式转移排序；不是未知图归纳评测。原管线对齐和可评价过滤后有 88,197 个训练问题、88,073 个 evaluation 问题，主表不能称为全部官方问题成绩。十一领域等权汇总，不能当作十一独立 benchmark。

MetaQA 使用官方 3-hop vanilla、完整固定电影 KB 和原 train/dev/test\cite{metaqa}。训练包含 114,196 题，test 为 14,274 题。普通有向关系加入显式逆向边；原始 KB 为 134,741 个三元组。topic 来自问题中的括号标注，测试不因答案不可达而删题。本轮复用已有三策略的真实结果并增加正例重标记。

PathQuestion 使用作者发布的 2H/3H 问题和知识库\cite{pq}，合并后有 2,256 个实体、3,377 个去重三元组。按完全相同的问题/topic 分组并合并答案，以种子 20261003 做 80/10/10 划分，train/valid/test 为 5,684/710/711；这是声明的自定义划分。训练中 5,537 题可构造至多三跳答案路径。标注链只提取 topic，所选路径重新在固定图上计算，不把原标注链当候选。

KQA Pro 的完整问题包含属性、集合操作和数值约束\cite{kqa}。本轮选取终端操作为 \texttt{What}、且恰有一个无歧义 \texttt{Find} topic 的实体输出问题，并将 KB 投影为实体关系图。train 子集随机划出 10\% 作为独立 valid，得到 3,205/357 题；官方 val 子集的全部 468 题为离线 evaluation。官方 test 不公开所需答案及程序，因此没有本轮官方 test 成绩。固定图有 16,960 个实体、177,923 个三元组，属性、概念和限定符不进入图。发布的 topic 注释可用，其余程序结构不进入检索；这不是完整 KQA Pro 组合推理任务。

我们还实际下载检查了 WebQSP 问题标注。现成 GraftNet Freebase 子图的发布说明包含基于问题 semantic parse 的关系筛选，因此本轮不使用它作为独立检索图。GrailQA 的完整数据和 Freebase 服务未就绪。暂缓这两者来自图来源和接口边界，不来自试跑后的指标；本研究不以知名基准数量代替对监督冲突的有效检验。

\subsection{控制变量与训练协议}
种子为 42--46；同一数据和种子下四策略读相同 prepared tensor、表示、原标签分层划分，共享初始化、batch 顺序、Adam 学习率 $10^{-4}$、最多 50 epoch、patience 10、最小改善 $10^{-5}$。WikiTopics batch 为 32，其余为 512。早停使用训练候选的 80/20 内部划分，重标记仍按原标签划分；没有根据测试表现选择策略、候选宽度或阈值。

独立 valid 文件保存供复现，本轮没有把它们用于策略专属调参。候选级早停允许同一问题的候选分布在两侧，不能替代问题级泛化验证。既有结果与新增对照沿用这一协议。GPU 上使用物化或按索引取相同 FP32 特征以适应显存，不改变特征计算、batch 或优化步骤。新数据统一使用训练入口，旧结果按其完整五种子来源复用。

\subsection{路径完成指标与候选上限}
令 $\operatorname{rank}_q(t)$ 为转移排名，评价所有固定候选可组成的有向答案路径，而不要求其为训练最短路。首次完整答案路径的前缀长度为
\begin{equation}
r_q=\min_{a\in A_q}\ \min_{P\subseteq C_q:\ s_q\leadsto a}\ \max_{t\in P}\operatorname{rank}_q(t).
\end{equation}
不存在完整路径时 $r_q=\infty$。APC-MRR 为 $\operatorname{mean}_q(1/r_q)$，Reach@$K$ 为 $\operatorname{mean}_q\mathbf1[r_q\le K]$；topic 本身为答案时按实现约定取 $r_q=1$。Candidate Oracle 表示在全部固定候选中存在至少一条完整答案路径。它区分候选覆盖与模型排序，也说明为什么不可达题无法仅靠重排恢复。

我们保留主分母，再报告同一候选池的 Oracle 可达条件指标。普通 KG 的选定 evaluation split 中无候选、图外答案或不可达题均记零。APC-MRR 是转移排序的路径完成指标，不是标准端到端 KGQA 的答案准确率。

\subsection{配对统计}
先对每个问题平均五个配对种子，再对 PCN Mask 与三个控制的差值做 10,000 次题级有放回 bootstrap，报告 95\% 百分位区间，bootstrap seed 为 20261003。WikiTopics 在领域内重采样后等权平均；其他数据按问题平均。多个指标共享抽样，重复差值用多项分布计数压缩，与从经验题分布采样等价。区间条件于这五个训练运行；不是全部训练随机性的总体区间，也不作多重比较校正。

\section{实际 PCN 的跨数据分布：RQ1}
表~\ref{tab:prevalence} 先报告问题是否存在。负例和 PCN 总数对五种子累加；题数按不同问题计数，受影响题数指至少一个种子出现 PCN。不能用这项比例代表单次运行的受影响概率。WikiTopics 在十一领域总计 53,074/88,197 题受影响，题级微平均为 60.18\%；它与领域等权的比例不同。

@@prevalence_table@@

PCN 并非只在稠密超图中出现。MetaQA 普通关系图比例为 11.81\%，高于 WikiTopics；PathQuestion 为 8.80\%，KQAPro 实体子集为 1.45\%。仅有多条最短路也不足以推断 PCN 数量：PathQuestion 的 2,249 个训练问题和 KQAPro 的 1,727 个训练问题具有多条最短路，但其实际 PCN 受路径引导负例池和采样数量约束。

逐题分布具有不均匀性。MetaQA seed 42 的受影响题为 71,753，平均 PCN 数为 4.659，中位数为 2；五种子任一次受影响题为 95,347。对应的直方图及每种子统计随结果提供。这个差异提醒我们，汇总比例既受图中路径结构影响，也受本项目的实际采样器影响；它不是整个知识图谱中语义伪负例的比例。

\section{共享候选上的排序结果：RQ2--RQ4}
\subsection{四策略主结果}
表~\ref{tab:main} 报告五种子均值，表~\ref{tab:delta} 给出配对差值与区间。所有数值以百分数或百分点表示。WikiTopics 为领域等权；其余按各自声明的 evaluation 问题平均。MetaQA 旧三策略逐题 CSV 的 Reach@5 由 $\mathrm{RR}\ge1/5$ 精确还原，并与已有每种子 JSON 指标核对，不重新训练或评价。

@@main_table@@
@@comparisons_table@@

PCN Mask 相对 Original 的 APC-MRR 改善分别为 WikiTopics 0.257、MetaQA 1.323、PathQuestion 2.179 和 KQAPro-entity 1.132 个百分点；四个配对区间均在零以上。相对 Random Mask 的改善分别为 0.235、1.528、1.449 和 0.741 个百分点。这些比较使用相同候选和逐题相同屏蔽数，因而支持处理位置与改善有关；它们尚未独立确认梯度变化的因果路径。

次指标的证据较不一致。PathQuestion 的 Reach@10 相对 Original 只提高 0.619 个百分点，区间下界近零，相对 Random Mask 的区间跨零；KQAPro 的 Reach@5 部分比较也跨零。更高 APC-MRR 并不要求每个 top-$K$ 指标都能被区间区分，不应将主指标结论复制到全部指标。

\begin{figure}[tbp]\centering
\begin{lrbox}{\pcnfigurebox}
@@evidence_pgf@@
\end{lrbox}\resizebox{\linewidth}{!}{\usebox{\pcnfigurebox}}
\caption{跨数据集现象及配对排序差异。左为真实五种子采样的 PCN 占比；右为 PCN Mask 减去 Original、Random Mask 和 Positive Relabel 的 APC-MRR 差值。点为逐题五种子均值差，误差线为 10,000 次题级配对 bootstrap 的 95\% 区间；WikiTopics 为领域内重采样后的等权宏平均。区间未作多重比较校正。}\label{fig:evidence}
\end{figure}

\subsection{撤销负监督与赋予正监督}
@@relabel_result_text@@

PathQuestion 上，PCN Mask 与 Positive Relabel 的 APC-MRR 差值为 0.295 个百分点，95\% CI 为 $[-0.302,0.875]$，当前样本无法明确区分。KQAPro 上重标记均值为 21.336\%，屏蔽为 21.113\%；屏蔽减重标记为 $-0.223$ 个百分点，区间上界极接近零。这不是屏蔽在所有数据上更好的证据，也不应把边界区间解释成稳健的普适反向优势。

重标记在部分数据上取得相当或更高均值，与结构一致转移可能提供有用正信号相符，但不能证明每个转移都语义正确。屏蔽是在缺乏确定标签依据时保留较弱判断的选择，其经验价值不能代替对重标记的检验。

\subsection{候选可达条件结果}
@@conditional_table@@

WikiTopics 的领域等权 Candidate Oracle 为 32.73\%，题级微平均为 35.34\%，MetaQA 为 56.02\%，PathQuestion 为 93.95\%，KQAPro-entity 为 82.69\%。这说明跨数据绝对 APC-MRR 差异不仅来自评分模型，也来自候选覆盖率及任务分布。

在 Oracle 可达题上，PathQuestion 的 Original/PCN Mask 条件 APC-MRR 为 45.68\%/48.00\%，KQAPro 为 24.16\%/25.53\%。不可达题的指标恒为零，故该条件分析更直接描述已有路径的排序；它不替代完整分母。WikiTopics 必须先逐领域求条件均值再宏平均，不能直接用宏平均 APC-MRR 除以宏平均 Oracle。

\subsection{全最短路径构造对照：RQ6}
@@all_shortest_table@@
@@all_shortest_result_text@@

这一对照在构造阶段改变正例与新负例，与四策略主表的“只改变监督处理”问题不同。即使它取得较高得分，也不能据此判定 PCN 身份足以支持任意语义正标签；如果它不占优，也不能否定其他多路径学习目标。当前实现仅覆盖普通 KG 的两个新基准，未替换 WikiTopics 的既有超图监督构造。

\section{何时有效：结构与排名的描述性分析}
\subsection{路径数、跳数与关系多样性}
在 evaluation 排名完成后，使用答案计算完整最短路径结构并分桶。训练监督和检索均不读取这些分析标签。图~\ref{fig:mechanism} 与源数据保留全部题，零最短路径桶独立记录。

PathQuestion 一跳桶的 270 题平均 APC-MRR 改善为 5.079 个百分点，二跳 239 题为 0.405，三跳 183 题为 0.445。按完整最短路径总数分组，1 条路径的 385 题改善 0.719 个百分点，2--3 条路径的 243 题改善 5.002，4--15 条路径的 63 题改善 0.903。路径更多或更长并不对应单调更大收益；多答案数和具体关系也会同时变化。

KQAPro 的一、二、三跳桶分别为 305、96、66 题，改善 1.134、0.666、1.817 个百分点。所选路径关系类型数 1、2--3、4--7 的桶分别改善 1.057、0.530、2.313 个百分点。这些分布与 PathQuestion 不同，支持继续考察问题分布；小桶不足以作机制定论。

\subsection{实际训练 PCN 密度与候选限制}
为了关联\emph{实际进入训练损失}的 PCN 和 held-out 收益，我们按 topic 聚合训练问题的 PCN/负例及平均每题 PCN，再匹配 evaluation topic。PathQuestion 有 691/711 题匹配，其中一题对应 topic 没有可用负例，密度分桶包含 690 题；KQAPro 匹配 190/468 题。未匹配或比例未定义的题不被置为零，也不进入对应密度桶。该统计不是利用 evaluation 答案重新采样得到的假想 PCN。

PathQuestion 的 topic PCN 比例从 $[0,0.01)$、$[0.01,0.05)$、$[0.05,0.10)$、$[0.10,0.25)$ 到 $[0.25,1]$ 时，对应平均 APC-MRR 改善为 0.623、1.374、2.352、3.757、4.782 个百分点，题数分别为 244、91、121、158、76。这一匹配子集呈递增趋势；不能据此把 topic 密度当作独立因果因素。KQAPro 高密度桶只有 2--7 题，无法支持同等趋势判断。

候选池大小和 Oracle 另作分桶。不可达题的改善恒为零是指标定义的结果，不是模型学到了排除失败。训练 topic 密度、最短路多重性、答案数、关系和候选覆盖存在共同变化，本文不拟合一个解释所有收益的评分表，也不以相关性推导策略选择规则。

\begin{figure}[tbp]\centering
\begin{lrbox}{\pcnfigurebox}
@@mechanism_pgf@@
\end{lrbox}\resizebox{\linewidth}{!}{\usebox{\pcnfigurebox}}
\caption{新数据上的描述性分桶。左：匹配 topic 的实际训练 PCN 密度与 evaluation APC-MRR 改善；点旁为题数，KQAPro 仅 190 题匹配且高密度桶很小。右：不同最大最短答案跳数的改善，点旁为题数。中心为五种子逐题均值的桶平均，无逐桶显著性或因果解释。所有分桶定义及未匹配题保留在源数据中。}\label{fig:mechanism}
\end{figure}

\subsection{领域与实际 PCN 的最终排名}
@@domain_result_text@@

在实际训练候选池中，五种子的 5,048 个 PathQuestion PCN 在 Original/PCN Mask 下平均 rank 为 3.707/3.379，中位数均为 3；KQAPro 的 546 个 PCN 为 3.676/3.595，中位数均为 2。top-10 比例分别从 98.02\% 到 98.16\%、94.69\% 到 95.05\%。这些训练池的 rank 变化描述被处理对象的行为，不能作为 held-out 泛化改善的独立证据。诊断中的分数并列按训练候选输入顺序稳定处理。

\section{相关工作}
\subsection{路径弱监督与多路径推理}
知识图谱问答可从问题/答案学习关系链，MetaQA 的变分推理和 IRN 的逐跳网络分别提供了相应任务与模型\cite{metaqa,pq}。KQA Pro 则用显式程序刻画组合推理\cite{kqa}。这些工作研究如何形成或利用推理链；本文以现有路径监督为对象，考察选择一条路径后实际负采样是否仍包含另一条完整最短答案路径。全最短路对照也不是对所有多路径推理算法的比较。

\subsection{难负例与伪负例}
DPR 将负例用于密集检索训练，ANCE 使用近邻负例改善训练难度\cite{dpr,ance}。难度高不意味着标签可靠。Debiased Contrastive Learning 对负采样可能包含同类别样本进行目标修正\cite{debiased}，密集检索中的置信度正则化也直接处理伪负例影响\cite{false_negative}。本文采用不同证据来源：在实际采样集合上检查完整图路径条件，不通过模型相似度或置信度推断语义正负。PCN 是结构限定的监督冲突子集，不是所有语义伪负例的替代定义。

\subsection{鲁棒标签学习与监督弃权}
噪声鲁棒交叉熵研究标签错误下的目标函数\cite{gce}，Snorkel 将弱监督来源及其冲突建模为训练信号\cite{snorkel}。本文没有建立通用标签噪声模型，而是定位一种可由图结构精确识别的冲突，再撤销对应负项。弃权的含义是这一监督来源不再对该候选作确定负判断，不能理解为候选被从检索图删除，或其语义已被自动校正。

\subsection{图与超图检索}
HyperRAG 将高阶事实与关系链检索结合\cite{hyperrag}。本文保留其发布评分器及 WikiTopics 表示，同时在普通 KG 上检验同一完整路径判据，研究范围是监督一致性而不是新增图表示或搜索器。结构形式从代价 1 的关系边变为代价 2 的二部转移时，PCN 的“实际负例与完整最短答案路径交集”含义保持不变；候选流程与问题分布仍须分别说明。

\section{讨论、局限与结论}
\subsection{结构证据的监督边界}
最短路径是一种弱监督启发式，而非语义真值。路径一致性能够指出未被选择这一理由不足以支持确定负标签，却不能确认关系链回答了问题。PCN Mask 对这一不确定性弃权；Positive Relabel 作更强正判断。重标记的非统一结果以及全最短路构造在两个新基准上的优势，均排除了将最小修正写为普适最优方案。后者重建正例覆盖和负例池，解决的问题比原集合内的冲突撤销更广；其额外收益应作为主张边界陈述，不能只放在可选实验附注。

等量随机对照控制了屏蔽数量，不能完全排除其他优化解释，例如屏蔽位置改变了 batch 的有效归一化或局部梯度分布。路径条件提供了明确定义且可复现的处理位置；分桶分析与训练 rank 尚不足以证明一个唯一因果机制。更强因果检验需要重新设计干预而非重复解释同一相关性。

\subsection{适用范围与局限}
证据覆盖有 topic 标注、固定语料和三跳构造的转移排序，以及一个冻结表示的两层 MLP。PathQuestion 是模板问题和自定义随机划分；KQAPro 只包含实体单 topic 子集并投影掉属性/限定符，evaluation 是官方 val 子集；WikiTopics 复用既有可评价分母和传导式图。本文不能声称完成全量 KQA Pro、实体链接、开放世界 KGQA 或端到端 RAG 的性能验证。

候选覆盖为方法设定上限，尤其 WikiTopics 和 MetaQA 的 Oracle 明显低于 100\%。本方法不扩大候选，无法补回丢失的完整路径。最短路定义也遗漏更长但语义正确的证据；关系标签和答案标注本身可能不完备。五种子的题级 CI 衡量这组模型在题分布上的差异，训练随机性、不同行为模型及多重比较仍限制结论。

\subsection{结论}
选择一条最短答案路径后，其他完整最短路径上的转移仍可能进入实际负损失。PCN 将该冲突定义为可核对的采样交集，完整 prefix--transition--suffix 条件排除局部距离下降的错误判据。保持标签、撤销负项在四组界定的数据上改善 APC-MRR，且高于逐题等量随机屏蔽；领域内仍有负结果。重标记没有统一劣势，全最短路构造在两个新基准上更高，因此最小撤销的贡献是保留既有监督集合时的有效修正，不是取代多路径监督。结构冲突识别与语义正标签赋予应分别论证，并在共同候选和明确的监督构造范围内评价。

\appendix
\section{复现、正确性与证据来源}
核心测试覆盖唯一最短路径、两条等长答案路径、绕远 prefix 的局部距离反例、屏蔽不改标签、仅 PCN 重标记、逐题随机数量、推理不读 PCN、替换 evaluation 答案不影响候选和特征。还验证了固定 batch 顺序下 indexed 与 materialized 特征的一 epoch CPU checkpoint 等价。训练配置、逐题评价、五种子汇总、分桶源数据和作图脚本随仓库提供；已完成旧结果保持原目录，失败任务的日志也保留。

本稿的表和图由同一真实聚合文件生成。主表不读取开发指标，缺失五种子时明确写 TODO。可选全最短路只在完成五种子后写数值。按索引取特征和同卡并发用于六张 RTX 3090 的吞吐，不修改模型、batch 或精度；数据获取失败时经本地文件中转，而不是生成替代数字。

\section{WikiTopics 中的真实结构案例}
已保存的训练案例问“哪些作曲家与 Giancarlo Santi 在同一类型中工作”。所选路径经电影事实超边到达 Luis Bacalov 等答案；实际采样负转移从 Giancarlo Santi 的超边到达 Spaghetti Western，而另一条完整答案路径继续经该类型相关的作曲家超边到达答案。二部距离满足 $0+2+2=4$。这个案例只证明固定超图中的完整路径资格；电影类型连接是否精确满足自然语言语义仍属语义问题。图~\ref{fig:overview} 的构造示例与此真实案例分开陈述，避免把插图当作数据计数证据。

\begin{thebibliography}{99}
\bibitem{hyperrag} W.-S. Lien, Y.-K. Chan, H.-L. Hsiao, et al. HyperRAG: Reasoning N-ary Facts over Hypergraphs for Retrieval Augmented Generation. arXiv:2602.14470, 2026. \url{https://arxiv.org/abs/2602.14470}.
\bibitem{metaqa} Y. Zhang, H. Dai, Z. Kozareva, A. Smola, L. Song. Variational Reasoning for Question Answering With Knowledge Graph. AAAI, 32(1), 2018. \url{https://doi.org/10.1609/aaai.v32i1.12057}.
\bibitem{pq} M. Zhou, M. Huang, X. Zhu. An Interpretable Reasoning Network for Multi-Relation Question Answering. COLING, pp. 2010--2022, 2018. \url{https://aclanthology.org/C18-1171/}.
\bibitem{kqa} S. Cao, J. Shi, L. Pan, et al. KQA Pro: A Dataset with Explicit Compositional Programs for Complex Question Answering over Knowledge Base. ACL, pp. 6101--6119, 2022. \url{https://doi.org/10.18653/v1/2022.acl-long.422}.
\bibitem{dpr} V. Karpukhin, B. Oguz, S. Min, et al. Dense Passage Retrieval for Open-Domain Question Answering. EMNLP, pp. 6769--6781, 2020. \url{https://doi.org/10.18653/v1/2020.emnlp-main.550}.
\bibitem{ance} L. Xiong, C. Xiong, Y. Li, et al. Approximate Nearest Neighbor Negative Contrastive Learning for Dense Text Retrieval. arXiv:2007.00808, 2020. \url{https://arxiv.org/abs/2007.00808}.
\bibitem{debiased} C.-Y. Chuang, J. W. Robinson, Y.-C. Lin, A. Torralba, S. Jegelka. Debiased Contrastive Learning. NeurIPS 33, 2020. \url{https://proceedings.neurips.cc/paper/2020/hash/63c3ddcc7b23daa1e42dc41f9a44a873-Abstract.html}.
\bibitem{false_negative} S. Wang, Y. Zhang, C.-T. Nguyen. Mitigating the Impact of False Negative in Dense Retrieval with Contrastive Confidence Regularization. AAAI, 38(17):19171--19179, 2024. \url{https://doi.org/10.1609/aaai.v38i17.29885}.
\bibitem{gce} Z. Zhang, M. R. Sabuncu. Generalized Cross Entropy Loss for Training Deep Neural Networks with Noisy Labels. NeurIPS 31, 2018. \url{https://proceedings.neurips.cc/paper/2018/hash/f2925f97bc13ad2852a7a551802feea0-Abstract.html}.
\bibitem{snorkel} A. Ratner, S. H. Bach, H. Ehrenberg, et al. Snorkel: Rapid Training Data Creation with Weak Supervision. Proceedings of the VLDB Endowment, 11(3):269--282, 2017. \url{https://www.vldb.org/pvldb/vol11/p269-ratner.pdf}.
\end{thebibliography}
\end{document}
